% File: Volume-II-Applications/appendices/appG-statistical-methods.tex

\chapter{Statistical Methods for Cosmology}
\label{app:cosmo-stats}

This appendix summarizes statistical techniques used to confront
theoretical cosmology with observational data, with emphasis on their
application to the lamphrodynamic dynamics of \rsvpabbr{}.  A rigorous
treatment of full Bayesian inference, likelihood functions, and
model comparison is necessary for the quantitative tests developed in
Chapters~8--10.  The following review parallels standard references
(e.g.\ Dodelson, Trotta) but adapts the presentation to the specific
parametrization of \rsvpabbr{}.

\section{Parameters and Priors}

Let $\theta$ denote a vector of cosmological parameters,
encompassing both standard parameters (e.g.\ present--day entropy
density $S_0$, Hubble parameter $H_0$, baryon fraction) and
lamphrodynamic parameters (e.g.\ coupling coefficients appearing in
the lamphrodynamic field equations).  The posterior distribution over
parameters given data $D$ is
\begin{equation}
\label{eq:bayes}
P(\theta\mid D) \propto \mathcal L(D\mid \theta)\,\pi(\theta),
\end{equation}
where $\mathcal L$ is the likelihood and $\pi$ is the prior.  Choice of
prior $\pi(\theta)$ should be explicitly justified; standard practice
includes uniform, Jeffreys, or weakly informative priors, depending on
physical interpretation and parameter scaling.

\section{Likelihoods}

For Gaussian--distributed data with covariance $\mathsf C(\theta)$ and
mean $\mu(\theta)$,
\begin{equation}
\mathcal L(D\mid \theta)
\propto |\mathsf C(\theta)|^{-1/2}
\exp\Bigl[-\tfrac12 (D-\mu(\theta))^{\sf T}\mathsf C(\theta)^{-1}
(D-\mu(\theta))\Bigr],
\end{equation}
where $D$ may include CMB power spectra $C_\ell$, matter power
spectra $P(k)$, and Type~Ia supernova distance moduli $\mu(z)$.
In \rsvpabbr{}, the dependence of $\mu(\theta)$ and $\mathsf C(\theta)$
on lamphrodynamic parameters enters through modified growth functions
and redshift relations (see Chapters~5--7).

\section{Maximum Likelihood and $\chi^2$ Statistics}

For independent Gaussian errors, maximizing $\mathcal L$ is equivalent
to minimizing $\chi^2$:
\begin{equation}
\chi^2(\theta) = \sum_i \frac{\bigl(D_i -
\mu_i(\theta)\bigr)^2}{\sigma_i^2}.
\end{equation}
Confidence intervals follow from the usual $\Delta\chi^2$ prescription,
provided regularity conditions are satisfied.  In practice, the
covariance matrix is often non--diagonal, requiring full matrix
estimators and noise modeling.

\section{Bayesian Evidence and Model Comparison}

Let $\mathcal M$ denote a model and $\mathcal E$ its Bayesian evidence,
\begin{equation}
\mathcal E(\mathcal M) =
\int \mathcal L(D\mid \theta,\mathcal M)\,\pi(\theta\mid\mathcal M)\,
\dd\theta.
\end{equation}
Model comparison proceeds by Bayes factors
\begin{equation}
B_{12} = \frac{\mathcal E(\mathcal M_1)}{\mathcal E(\mathcal M_2)},
\end{equation}
with conventional interpretive scales (e.g.\ Jeffreys’ scale).  Evidence
is especially important for comparing \rsvpabbr{} with $\Lambda$CDM,
because additional parameters (entropy--vector couplings) alter the
dimensionality of parameter space.

\section{Information Criteria}

Information criteria provide approximate model comparison without
computing full Bayesian evidence.  The Akaike information criterion is
\begin{equation}
\mathrm{AIC}=2k-2\ln\hat{\mathcal L},
\end{equation}
where $k$ is the number of parameters and $\hat{\mathcal L}$ is the
maximum likelihood.  The Bayesian information criterion is
\begin{equation}
\mathrm{BIC}=k\ln N - 2\ln\hat{\mathcal L},
\end{equation}
where $N$ is the number of data points.  Lower values indicate
preferred models, penalizing parameter complexity.

\section{Sampling Methods}

Sampling from the posterior \eqref{eq:bayes} typically requires
Markov--chain Monte Carlo (MCMC) or nested sampling algorithms.
Standard approaches (Metropolis--Hastings, Hamiltonian Monte Carlo,
\texttt{emcee}, \texttt{MultiNest}) can be adapted to \rsvpabbr{} by
supplying numerically integrated transfer functions and redshift
relations.  Parameter covariance can be evaluated by examining chain
autocorrelation times and convergence diagnostics (e.g.\ Gelman--Rubin
statistics).

\section{Forecasting and Fisher Information}

For parameter forecasting (e.g.\ next--generation CMB or large--scale
structure surveys), the Fisher matrix is defined by
\begin{equation}
F_{\alpha\beta}=
-\left\langle
\frac{\partial^2 \ln\mathcal L}{\partial\theta_\alpha
\partial\theta_\beta}
\right\rangle.
\end{equation}
Parameter uncertainties are approximated by
\[
\sigma(\theta_\alpha)\approx \sqrt{(F^{-1})_{\alpha\alpha}}.
\]
Fisher forecasting allows a quick estimate of the extent to which
future surveys can constrain lamphrodynamic parameters.

\section{Goodness of Fit}

Standard tests include reduced $\chi^2$,
Kolmogorov--Smirnov statistics, residual analysis, and cross--validation
of independent data sets (e.g.\ CMB and BAO).  Given the extended
parameter space of \rsvpabbr{}, it is prudent to examine potential
degeneracies between lamphrodynamic and standard cosmological
parameters (e.g.\ $H_0$, $\Omega_m$, spectral indices).

\section{Remarks and Further Reading}

Comprehensive introductions include Trotta (\emph{Bayes in Cosmology})
and Dodelson (\emph{Modern Cosmology}).  Recent data analysis
pipelines for CMB, BAO, and supernovae provide practical examples of
implementing parameter inference in extended cosmological models.
Adaptations to \rsvpabbr{} involve modified theoretical predictions
for $\mu(\theta)$ and $\mathsf C(\theta)$ and are discussed throughout
Chapters~8--10.

